
\section{Methods}
\label{sec:methods}

\textbf{Pretraining} We pretrained causal language models \citep{vaswani2017attention, brown2020language} using the Qwen 3 architecture \citep{yang2025qwen3technicalreport}, sweeping model sizes: $34$M, $62$M, $93$M, $153$M, $344$M.
For each model size, we created multiple pretraining corpora by contaminating a high quality corpus \citep{penedo2024finewebdatasetsdecantingweb} with a different number of replicas of the benchmark test set: from $0$ (uncontaminated) through $1, 3, 10, 32, 100, 316, 1000, 3162$ (logarithmically uniform).
Pretraining floating point operations (FLOP) $C$ was approximated $ 6  \, N \, D$ \citep{kaplan2020scaling,sardana2024beyond,porian2024resolving,gadre2024languagemodelsscalereliably}, where $N$ is model parameters and $D$ is tokens.
Each model was pretrained with a target of 20 tokens-per-parameter \citep{hoffman2022chinchilla}; the realized budget was uniformly $71.4\%$ of this target (App.~\ref{app:sec:pretraining_implementation_details}), which rescales all runs by a common factor and affects no comparison.
For each tuple of parameters and test set replicas, we trained one model.
For details, see App.~\ref{app:sec:pretraining_implementation_details}.

As a key distinction, \citet{bordt2025howmuch} pretrains models up to ~1.6B, but achieves this by contaminating disjoint subsets of the test set at different replica counts within a single training run — a cheaper but messier design. We pretrain separate models per contamination amount for tighter experimental control, which is significantly more expensive and thus limits our maximum scale.

\textbf{Benchmark} We chose the MATH \citep{hendrycks2021measuringmathematicalproblemsolving} benchmark of competition math problems because it is comparatively large (5,000 test problems), and includes solutions as well as verifiable answers. 
The MATH test set contains $\mathord{\sim}1.4$M tokens under the Qwen 3 tokenizer.

\textbf{Evaluation} 
The first metric we report is \emph{Math Verify} \citep{kydlicek2025fixing}, defined as the fraction of problems for which the model generates solutions verified to be mathematically equivalent to the boxed answers.
We initially used EleutherAI's LM Evaluation Harness \citep{gao2024evalharness}, but discovered a bug in how Math Verify is computed; for example, the benchmark's gold reference solutions obtained a score of $\mathord{\sim}70\%$.
We worked with the developers to correct the implementation (Appendix~\ref{app:sec:math_verify_bug}).
This suggests that \texttt{math\_minerva} scores using the Eval Harness prior to Aug 23, 2025 may be incorrect.
The second metric we report is the \emph{Cross Entropy} of the gold references, which was previously demonstrated to be useful for studying scaling properties of generative evaluations during pretraining \citep{schaeffer2025pretrainingscalinglawsgenerative}.
We used temperature-only sampling \citep{schaeffer2025minpmaxexaggerationcritical}, beginning with temperature $0$, and expanding to more temperatures in Sec.~\ref{sec:inference}.